\documentclass[11pt]{article}
\usepackage[a4paper,margin=27mm]{geometry}
\usepackage{amsmath,amssymb,amsthm,booktabs,array,microtype}
\usepackage[hidelinks]{hyperref}

\newtheorem{theorem}{Theorem}
\newtheorem{lemma}{Lemma}
\newcommand{\C}{\mathbb C}
\newcommand{\R}{\mathbb R}
\newcommand{\PP}{\mathbb P}
\newcommand{\one}{\mathbf 1}

\title{A Five-Point Counterexample to Concavity of Determinantal Entropy}
\author{}
\date{Review note, 5 September 2026}

\begin{document}
\maketitle

\begin{abstract}
Lyons asked whether the Shannon entropy of a finite determinantal probability
measure is concave in its positive-contraction kernel.  We give a counterexample
on five labelled points when complex Hermitian kernels are allowed.  The
construction comes from a rank-two projection and a purely imaginary direction
that is invisible to first order in every event probability.  The usually
negative Fisher-information term therefore vanishes, while the second-order
bending of the probability map has positive entropy pairing.  Four short integer
vectors specify the example, and a dependency-free exact checker proves
\(\tfrac12<D^2H_K[A,A]<\tfrac35\) at the explicit interior kernel.
\end{abstract}

\section{The question and the answer}

Let \(E=[5]\).  If \(K\) is a Hermitian matrix with \(0\preceq K\preceq I\),
the determinantal probability measure \(\PP^K\) on subsets of \(E\) is defined by
\[
 \PP^K(T\subseteq X)=\det K[T].
\]
Writing \(p_S(K)=\PP^K(X=S)\), its entropy is
\[
 H(K)=-\sum_{S\subseteq E}p_S(K)\log p_S(K).
\]
The conjecture that \(H\) is concave in \(K\) was posed by
Lyons~\cite{Lyons2003DPM} and restated as Conjecture~2.6 in his ICM
survey~\cite{Lyons2014}.  The formulation there permits real or complex Hilbert
spaces.

Here is the counterexample.  Put
\[
\begin{aligned}
 \one&=(1,1,1,1,1)^T,             &x&=(-2,-1,0,1,2)^T,\\
 q_2&=(2,-1,-2,-1,2)^T,           &q_3&=(-1,2,0,-2,1)^T,
\end{aligned}
\]
and write \(u\wedge v=uv^T-vu^T\).  Define
\begin{equation}\label{eq:construction}
 P=\frac15\one\one^T+\frac1{10}xx^T,
 \qquad
 B=\one\wedge x-\frac12q_2\wedge q_3,
 \qquad A=iB,
\end{equation}
and
\begin{equation}\label{eq:K}
 \varepsilon=10^{-6},\qquad K=\varepsilon I+(1-2\varepsilon)P.
\end{equation}

\begin{theorem}\label{thm:main}
The matrices in \eqref{eq:construction}--\eqref{eq:K} satisfy \(0<K<I\),
\(A=A^*\), and
\[
 \frac12<D^2H_K[A,A]<\frac35.
\]
Consequently \(H\) is not concave on the complex Hermitian positive
contractions on \(\C^5\).
\end{theorem}

The rest of the note explains both the mechanism and the exact verification.
The displayed decimal value, included only for orientation, is
\[
 D^2H_K[A,A]=0.56612746277851002589\ldots.
\]

\section{Why a purely imaginary direction is useful}

The complete-event probabilities admit the signed determinant formula
\begin{equation}\label{eq:event}
 p_S(K)=(-1)^{|E\setminus S|}
 \det\bigl(K-I_{E\setminus S}\bigr),
\end{equation}
where \(I_T\) is the diagonal coordinate projection onto \(T\).  This follows
directly from M\"obius inversion of the inclusion probabilities.

Let \(\Phi(K)=(p_S(K))_S\) be the map from kernels to the probability simplex.
At an interior kernel, differentiating entropy twice gives
\begin{equation}\label{eq:hessian}
 D^2H_K[A,A]
 =-\sum_S\frac{(Dp_S(K)[A])^2}{p_S(K)}
  -\sum_S D^2p_S(K)[A,A]\log p_S(K).
\end{equation}
The first sum is the pullback of the Fisher metric and is always non-positive.
The second sum records the curvature of the probability map in the direction of
the entropy gradient.

The construction makes the first sum disappear.  If \(K\) is real symmetric
and \(B\) is real skew-symmetric, then
\[
 \det(K-I_{E\setminus S}+itB)
 =\det(K-I_{E\setminus S}-itB),
\]
because transposition changes \(t\) to \(-t\).  Thus every event probability is
an even function of \(t\), and
\begin{equation}\label{eq:vertical}
 Dp_S(K)[iB]=0\qquad(S\subseteq E).
\end{equation}
If
\[
 b_S(K,B)=\left.\frac{d^2}{dt^2}p_S(K+itB)\right|_{t=0},
\]
then \eqref{eq:hessian} reduces to
\begin{equation}\label{eq:reduced}
 D^2H_K[iB,iB]= -\sum_S b_S(K,B)\log p_S(K).
\end{equation}
In information-geometric language, \(iB\) lies in the differential kernel of
\(\Phi\); the sign is decided entirely by the vertical second-order bending of
\(\Phi\).

\section{How the four vectors were chosen}

The vectors \(\one,x,q_2,q_3\) are pairwise orthogonal across the two displayed
pairs: \(\one\perp x\), and \(q_2,q_3\perp\one,x\).  Since
\(\|\one\|^2=5\) and \(\|x\|^2=10\), the matrix \(P\) is the orthogonal
projection onto \(\langle\one,x\rangle\).  The two summands of \(B\) rotate,
respectively, inside \(\mathrm{ran}(P)\) and inside \(\ker(P)\).  Hence
\begin{equation}\label{eq:commute}
 P^2=P=P^T,\qquad B^T=-B,\qquad PB=BP.
\end{equation}
It follows at once that \(K\) has eigenvalue \(1-\varepsilon\) on
\(\mathrm{ran}(P)\) and eigenvalue \(\varepsilon\) on \(\ker(P)\).  Thus
\(0<K<I\), while \(A=iB\) is Hermitian.

The two pieces of \(B\) are both needed.  A rotation confined to only one side
of the projection gives a logarithmic divergence in \eqref{eq:reduced}.  The
relative coefficient \(-\tfrac12\) balances the two sides so that the divergent
terms cancel, leaving a finite positive limit.

\section{The positive boundary limit}

For \(0<\delta<1/2\), set
\(K_\delta=\delta I+(1-2\delta)P\).  The coordinates of \(x\) are distinct,
so every two-row minor of the matrix with columns \(\one,x\) is nonzero.
Consequently the rank-two projection \(P\) is full spark.  For each subset \(S\)
there is a positive rational number \(a_S\) such that
\begin{equation}\label{eq:asymptotic}
 p_S(K_\delta)=a_S\delta^{w_S}(1+O(\delta)),
 \qquad w_S=\bigl||S|-2\bigr|.
\end{equation}
More explicitly,
\[
 a_S=\det P[S]\quad(|S|\le2),
 \qquad
 a_S=\det(I-P)[E\setminus S]\quad(|S|\ge2).
\]
These formulas follow either from the spectral Bernoulli mixture for
\(K_\delta\), or directly from principal-minor expansion and Jacobi's
complementary-minor identity.

Let \(b_S=b_S(P,B)\).  Direct second-order determinant expansion gives the two
exact cancellations
\begin{equation}\label{eq:cancellations}
 \sum_S b_S=0,
 \qquad
 \sum_S w_Sb_S=0.
\end{equation}
The first is forced by \(\sum_Sp_S=1\).  The second is the balancing condition
between the rotations in \(\mathrm{ran}(P)\) and \(\ker(P)\).  Equations
\eqref{eq:reduced}--\eqref{eq:cancellations}, together with
\(\delta\log\delta\to0\), now give
\begin{align}
 \lim_{\delta\downarrow0}D^2H_{K_\delta}[iB,iB]
 &=-\sum_Sb_S\log a_S \notag\\
 &=\frac15\log
 \frac{2^{260}3^{143}5^{350}}{7^{87}13^{284}}
 =0.57665861803438136898\ldots.       \label{eq:limit}
\end{align}
This is an exact positive number: the numerator inside the logarithm is a
392-digit integer, whereas the denominator is a 390-digit integer.

For transparency, Table~\ref{tab:grouped} groups the 32 determinant terms by
their value of \(a_S\).  Its second column is
\(B_a=\sum_{S:a_S=a}b_S\); hence \(-\sum_aB_a\log a\) reproduces
\eqref{eq:limit}.  The ungrouped table is supplied with the checker.

\begin{table}[ht]
\centering
\small
\caption{Grouped data for the boundary logarithmic sum.}\label{tab:grouped}
\begin{tabular}{@{}rr@{\qquad}rr@{}}
\toprule
\(a\)&\(B_a\)&\(a\)&\(B_a\)\\
\midrule
\(1/50\)&\(72/5\)&\(2/5\)&\(-2\)\\
\(2/25\)&\(-204/5\)&\(12/25\)&\(128/5\)\\
\(3/25\)&\(41/5\)&\(13/25\)&\(284/5\)\\
\(9/50\)&\(-606/5\)&\(3/5\)&\(120\)\\
\(1/5\)&\(20\)&\(7/10\)&\(-32\)\\
\(7/25\)&\(247/5\)&\(4/5\)&\(-36\)\\
\(3/10\)&\(60\)&\(1\)&\(-100\)\\
\(8/25\)&\(-112/5\)&&\\
\bottomrule
\end{tabular}
\end{table}

\section{An exact interior certificate}

The positive limit already proves that all sufficiently small \(\delta\) are
counterexamples.  To remove the hidden threshold, Theorem~\ref{thm:main} fixes
\(\varepsilon=10^{-6}\).  The following verification uses rational arithmetic
throughout.

For a positive rational \(r\in[1,2]\), set \(y=(r-1)/(r+1)\).  The positive
series for \(2\,\mathrm{atanh}(y)\) gives, for every positive integer \(m\),
\begin{equation}\label{eq:logbound}
 2\sum_{j=0}^{m-1}\frac{y^{2j+1}}{2j+1}
 \le \log r\le
 2\sum_{j=0}^{m-1}\frac{y^{2j+1}}{2j+1}
 +\frac{2y^{2m+1}}{(2m+1)(1-y^2)}.
\end{equation}
Every positive rational probability is written as \(2^kr\) with
\(1\le r<2\).  Taking \(m=24\) in \eqref{eq:logbound} gives directed rational
intervals for all logarithms in \eqref{eq:reduced}.

The event probabilities and second derivatives are also rational.  They are
computed by expanding each five-by-five determinant over its 120 permutations
and extracting the coefficients of \(t^0\) and \(t^2\).  Summation checks
\[
 p_S(K)>0\quad\hbox{for all }S,
 \qquad \sum_Sp_S(K)=1,
 \qquad \sum_Sb_S(K,B)=0.
\]
Applying the lower logarithm endpoint to positive coefficients and the upper
endpoint to negative coefficients, with the choices reversed for the upper
bound, proves
\begin{equation}\label{eq:exactinterval}
 0.5661274627785100258893
 <D^2H_K[iB,iB]<
 0.5661274627785100258923.
\end{equation}
The endpoints in the computation are exact fractions; the decimals in
\eqref{eq:exactinterval} are shortened outward displays.  In particular,
\(\tfrac12<D^2H_K[iB,iB]<\tfrac35\).

Because \(K\) is an interior positive contraction, \(K+tA\) remains a positive
contraction for all sufficiently small real \(t\).  A concave twice
differentiable function cannot have positive second derivative along such an
interior line.  This completes the proof of Theorem~\ref{thm:main}.

The exact checker uses only the Python standard library.  It is available with
the ungrouped boundary table at
\begin{center}
\small
\url{https://github.com/randomcat4/icm-conjecture-results/tree/release/I05-lyons-counterexample-20260905/results/I05_lyons_counterexample/author_review_note}.
\end{center}

\section{Scope}

The midpoint \(K\) is real symmetric, but the direction \(A=iB\) is purely
imaginary.  The example therefore disproves the conjecture with the real-or-
complex quantifier used in the cited sources.  It does not decide whether
entropy is concave when kernels and directions are required to be real
symmetric, and it does not establish that five is the smallest possible
dimension.

The calculation also suggests a broader structural separation near projection
kernels: the Fisher term can enforce negative curvature in all real directions
while commuting imaginary directions retain a finite positive boundary limit.
A companion note in the repository proves such a separation for an explicit
one-parameter family.  That extension is not needed for the counterexample
above.

We make no absolute priority claim.  To the best of our knowledge, the cited
conjecture had no previously published proof or counterexample; private
communications, seminars, unpublished notes, and unsuccessful attempts are
outside the scope of our search.

\bibliographystyle{plain}
\bibliography{references}

\end{document}
